\section{Objective, source audit, and the complexity target}
\label{sec:context}

The computational objective is to move the existing exact recognizer toward
an algorithm whose running time is bounded by
\(n^{(\log n)^{O(1)}}\), where \(n\) measures the input knot diagram. A useful
continuation must contribute an implemented step, a proved structural
bound, or a precisely exposed obstruction. Merely attaching a polynomial
subroutine to an exponential search does not settle this objective.

The present work targets compressed normal-surface analysis. A normal
surface can have exponentially many elementary disks while its coordinates
have only polynomially many bits. Thus a program that follows every normal
disk, every edge point, or every connected component explicitly is unsuitable
for the intended complexity target. The interval-pairing representation
already used by the repository is an appropriate starting point: its input
size is the number of pairings and the bit length of their endpoints.

\subsection{Pinned repository state}

All changes are relative to ProveIt commit
\begin{center}
\baselinecommit.
\end{center}
The \path{Topology/UnknotRecognition/fast} Git tree at that commit is
\begin{center}
\texttt{72834388419f5016d6db28ca0f028794f30b68db}.
\end{center}
The source was accessed through the selected GitHub connection. Locally
available copies were reused only after matching their Git blob contents to
the pinned tree; changed source files were fetched at the commit itself.
This distinction matters because the repository was receiving many concurrent
research integrations during the audit.

The maintained source contains substantially more than the original
Khovanov scanner. It includes compressed words and group certificates, braid
reductions, interval-component sharing, exact homological windows,
compressed twists, several Jones evaluators, and normal-surface topology
queries. We did not treat an old delivery snapshot as the current baseline.
The audit focused on
\path{fastunknot/interval_orbits.py},
\path{fastunknot/interval_orbit_verify.py},
\path{fastunknot/normal_surface_geometry.py}, and
\path{fastunknot/normal_surface_orbits.py}.
The corresponding maintained synthesis sections already explain interval
orbits, independent orbit proofs, coordinate multiplicity, and boundary
component classification \cite{ProveIt2026}.

Two limitations stand out in that source. Its periodic closure restarts a
lexicographic all-pairs search after every merge, even though only one
surviving pairing has changed. Its normal-surface topology query computes
aggregate Euler and component data, but its
\code{compressing\_disk} flag concerns the entire supplied surface being
a disk. The weighted component extraction needed to detect a disk hidden
inside a disconnected vector is explicitly absent.

The new work addresses these two limitations. It does not repeat the recent
compressed singleton-elimination or raw-contraction work on group
presentations. It also does not alter the meaning of the old
\code{compressing\_disk} field: a separate component query supplies the
stronger, differently quantified predicate.

\subsection{What the current literature establishes}

Agol, Hass, and Thurston proved polynomial-time counting for interval
pseudogroups, extended the algorithm to additive vector weights, and
applied it to compressed normal-surface components
\cite{AHT2006}. That classical result is the mathematical foundation of
the implementation here. In particular, the existence of a polynomial
algorithm for obtaining topology from supplied binary coordinates is not
a new theorem of this report.

The quasi-polynomial announcement must be stated with its actual
parameters. Lackenby's March 2021 Oxford handout gives
\[
 2^{O((\log n)^3)}
 =n^{O((\log n)^2)},
\]
with hierarchy length \(L=O((\log n)^2)\) and a polynomial surface-complexity
parameter \cite{Lackenby2021}. Earlier publicity used the stronger expression
\(n^{O(\log n)}\). Both are quasi-polynomial, but they are not identical
bounds and should not be substituted for one another in an implementation
claim.

Lackenby's July 2026 preprint gives a new hierarchy approach to
incompressibility and unknot recognition. Section 9 explicitly leaves the
running time of some operations unspecified; Proposition 9.1 bounds
iterations by \(L(g+1)^L\). Section 10 refines the hierarchy length, and
Theorem 14.4 treats compressed cutting and parallelity information
\cite{Lackenby2026}. These developments are highly relevant specifications
for future work, but the inspected preprint does not itself assert a
quasi-polynomial runtime for the full algorithm. The current repository
already records that qualification correctly.

\subsection{The deliverable's main statements}

The first result is an exact simulation theorem: the new scheduler produces
the same ordered sequence of successful periodic reductions as the pinned
implementation. It follows that all later carrier choices, transmissions,
truncations, and complete certificates agree as well. This is stronger than
merely observing equal final orbit counts.

The second result is a work bound that is independent of the expanded
interval size. A pure queue needs at most \((k-1)^2\) merger tests in one
closure. The adaptive default adds at most \(\binom{k}{2}\) preliminary
tests. The proof identifies precisely what is cached and what is invalidated;
it does not count failed arithmetic operations as free.

The third result exhibits a family on which the improvement changes
the asymptotic cost of a \emph{complete} count. Both old and new algorithms
take five cycles, so expensive later phases do not conceal the improvement.
This is an implementation-specific lower bound, not a lower bound for the
abstract orbit-counting problem.

The fourth result is a source-bound component predicate. Given a validated
compact orientable triangulation with one torus boundary and an admissible
normal vector, the new query counts all disk components and all
compressing-disk components without listing the components individually.
It uses five coordinates regardless of the triangulation size. This is
particularly useful when a candidate normal vector includes irrelevant
vertex links, closed components, or many parallel copies.

Throughout, \(k\) denotes an interval-pairing count, \(N\) its encoded
universe size, \(t\) a tetrahedron count, and \(n\) a knot-diagram size.
Bounds in \(k\) or \(t\) do not silently become bounds in \(n\). Any passage
between these parameters requires the corresponding geometric construction.
